% Part: computability
% Chapter: recursive-functions
% Section: pr-functions

\documentclass[../../../include/open-logic-section]{subfiles}

\begin{document}

\olfileid{cmp}{rec}{prf}
\olsection{بنسټيزې بازګشتي تابعې}

راځئ يو ځل بيا وليکو چې له شته تابعو څخه د بنسټيز بازګښت او ترکيب
په کارولو نوې تابعې څنګه تعريفولے شو۔

\begin{defn}
  \ollabel{defn:primitive-recursion} فرض کړئ~$f$ يوه~$k$-ځاييزه
  تابع ($k\ge 1$) او~$g$ يوه~$(k+2)$-ځاييزه تابع ده۔ هغه تابع چې
  \emph{له~$f$ او~$g$ څخه په بنسټيز بازګښت تعريفېږي}،
  $(k+1)$-ځاييزه تابع~$h$ ده چې لاندې معادلې ئې تعريفوي:
  \begin{align*}
    h(x_0,\dots,x_{k-1},0) & =  f(x_0,\dots,x_{k-1}) \\
    h(x_0,\dots,x_{k-1},y+1) & = g(x_0,\dots,x_{k-1}, y, h(x_0,\dots,x_{k-1}, y))
  \end{align*}
\end{defn}

\begin{defn}
  \ollabel{defn:composition} فرض کړئ~$f$ يوه~$k$-ځاييزه تابع ده، او
  $g_0$، \dots،~$g_{k-1}$ هغه~$k$ تابعې دي چې ټولې~$n$-ځاييزې
  دي۔ هغه تابع چې \emph{له~$f$ او~$g_0$، \dots،~$g_{k-1}$ څخه په
    ترکيب تعريفېږي}، $n$-ځاييزه تابع~$h$ ده چې داسې تعريفېږي:
  \[
  h(x_0, \dots, x_{n-1}) =
  f(g_0(x_0, \dots, x_{n-1}), \dots, g_{k-1}(x_0, \dots, x_{n-1})).
  \]
\end{defn}

له~$\Succ$ او پروجکشن تابعو
\[
\Proj{n}{i}(x_0,\dots,x_{n-1}) = x_i,
\]
سره سره، د هر طبيعي عدد~$n$ او~$i < n$ دپاره تابع~$\Zero(x) = 0$
هم په بنسټيزو بازګشتي تابعو کښې شاملوو۔

\begin{defn}
  د بنسټيزو بازګشتي تابعو سټ له~$\Nat^n$ څخه~$\Nat$ ته د تابعو هغه
  سټ دے چې په لاندې مادو استقرايي ډول تعريفېږي:
  \begin{enumerate}
  \item $\Zero$ بنسټيزه بازګشتي ده۔
  \item $\Succ$ بنسټيزه بازګشتي ده۔
  \item هره پروجکشن تابع~$\Proj{n}{i}$ بنسټيزه بازګشتي ده۔
  \item که~$f$ يوه~$k$-ځاييزه بنسټيزه بازګشتي تابع وي او~$g_0$،
    \dots،~$g_{k-1}$، $n$-ځاييزې بنسټيزې بازګشتي تابعې وي، نو له
    $g_0$، \dots،~$g_{k-1}$ سره د~$f$ ترکيب بنسټيزه بازګشتي دے۔
  \item که~$f$ يوه~$k$-ځاييزه بنسټيزه بازګشتي تابع او~$g$ يوه
    $k+2$-ځاييزه بنسټيزه بازګشتي تابع وي، نو له~$f$ او~$g$ څخه په
    بنسټيز بازګښت تعريف شوې تابع بنسټيزه بازګشتي ده۔
\end{enumerate}
\end{defn}

\begin{explain}
په لنډه وينا، د بنسټيزو بازګشتي تابعو سټ هغه تر ټولو وړوکے سټ دے
چې~$\Zero$، $\Succ$ او پروجکشن تابعې~$\Proj{n}{j}$ لري، او د ترکيب
او بنسټيز بازګښت له مخې بند دے۔

د بنسټيزو بازګشتي تابعو سټ د «پړاوونو» له مخې هم بيانولے شو۔
$S_0$ دې د پيلوونکو تابعو، يعنې~$\Zero$، $\Succ$ او پروجکشنونو، سټ
وښيي۔ دا د~$0$ پړاو بنسټيزې بازګشتي تابعې دي۔ کله چې يو
پړاو~$S_i$ تعريف شو، $S_{i+1}$ دې د پخپله~$S_i$ د تابعو او هغو
ټولو تابعو سټ وي چې د همدې پړاو پر تابعو د ترکيب يا بنسټيز بازګښت
يوه بېلګه په کارولو لاس ته راځي۔ بيا
\[
S = \bigcup_{i \in \Nat} S_i
\]
د ټولو بنسټيزو بازګشتي تابعو سټ دے۔
\textbf{د ژباړې د نسخې سمون (OLCMP-005):} سرچينې راتلونکی پړاو
يوازې د مخکيني پړاو پر تابعو د يو عمل له پايلو جوړاوه او پخپله ئې
مخکينۍ تابعې نۀ ساتلې؛ دلته پړاوونه په څرګنده تزايدي کړل شول، څو
اتحاد د تعريف شوو تابعو له مخې بند وي۔
\end{explain}

راځئ تصديق کړو چې~$\Add$ بنسټيزه بازګشتي تابع ده۔

\begin{prop}
  د جمع تابع~$\Add(x,y) = x+y$ بنسټيزه بازګشتي ده۔
\end{prop}

\begin{proof}
د~$\Add$ يو بنسټيز بازګشتي تعريف مو لا له دوو تابعو~$f$ او~$g$ څخه
لرو چې د~\olref{defn:primitive-recursion} له شکل سره سمون خوري:
\begin{align*}
  \Add(x_0, 0) & = f(x_0) = x_0\\
  \Add(x_0, y+1) & = g(x_0, y, \Add(x_0, y)) =
  \Succ(\Add(x_0, y))
\end{align*}
نو~$\Add$ په دې شرط بنسټيزه بازګشتي ده چې~$f$ او~$g$ هم داسې وي۔
$f(x_0) = x_0 = \Proj{1}{0}(x_0)$ ده، او پروجکشن تابعې بنسټيزې
بازګشتي شمېرل کېږي، نو~$f$ بنسټيزه بازګشتي ده۔ تابع~$g$ هغه
درې‌ځاييزه تابع~$g(x_0, y, z)$ ده چې داسې تعريفېږي:
\[
g(x_0, y, z) = \Succ(z).
\]
دا لا نۀ ښيي چې~$g$ بنسټيزه بازګشتي ده، ځکه~$g$ او~$\Succ$ په کره
توګه يوه تابع نۀ دي: $\Succ$ يوځاييزه ده، خو~$g$ بايد درې‌ځاييزه
وي۔ مګر~$g$ په ترکيب په «رسمي» ډول داسې تعريفولے شو:
\[
g(x_0, y, z) = \Succ(\Proj{3}{2}(x_0, y, z))
\]
څرنګه چې~$\Succ$ او~$\Proj{3}{2}$ بنسټيزې بازګشتي تابعې شمېرل
کېږي، نو~$g$ هم داسې ده، ځکه چې له بنسټيزو بازګشتي تابعو څخه په
ترکيب تعريفېدے شي۔
\end{proof}

\begin{prop}
  \ollabel{prop:mult-pr}
  د ضرب تابع~$\Mult(x,y) = x \cdot y$ بنسټيزه بازګشتي ده۔
\end{prop}

\begin{proof}
  تمرين۔
\end{proof}

\begin{prob}
  \olref[cmp][rec][prf]{prop:mult-pr} په دې ښودلو ثابته کړئ چې د
  $\Mult$ بنسټيز بازګشتي تعريف د
  \olref[cmp][rec][prf]{defn:primitive-recursion} غوښتل شوي شکل ته
  راوستل کېدے شي، او اړوندې تابعې~$f$ او~$g$ بنسټيزې بازګشتي دي۔
\end{prob}

\begin{ex}
دا مو د بنسټيز بازګشتي تعريف تر ټولو لومړۍ بېلګه ده:
\begin{align*}
h(0) & =  1 \\
h(y+1) & =  2 \cdot h(y).
\end{align*}
دا تابع د~\olref{defn:primitive-recursion} غوښتل شوي شکل ته نۀ شي
راتلے، ځکه~$k=0$ دے۔ په تعريف کښې ثابتونه~$1$ او~$2$ هم راځي۔ د
لومړۍ ستونزې د هوارولو دپاره يو بې‌کاره آرګومېنټ ورزياتوو او
تابع~$h'$ تعريفوو:
\begin{align*}
h'(x_0, 0) & =  f(x_0) = 1 \\
h'(x_0, y+1) & =  g(x_0, y, h'(x_0, y)) = 2 \cdot h'(x_0, y).
\end{align*}
تابع~$f(x_0) = 1$ له~$\Succ$ او~$\Zero$ څخه په ترکيب تعريفېدے شي:
$f(x_0) = \Succ(\Zero(x_0))$۔ تابع~$g$ له~$g'(z) = 2 \cdot z$ او
پروجکشنونو څخه په ترکيب تعريفېدے شي:
\begin{align*}
g(x_0, y, z) & = g'(\Proj{3}{2}(x_0, y, z))
\intertext{او~$g'$ بيا پخپله په ترکيب داسې تعريفېدے شي:}
g'(z) & = \Mult(g''(z), \Proj{1}{0}(z))
\intertext{او}
g''(z) & = \Succ(f(z)),
\end{align*}
چې~$f$ هماغه پورته تابع ده: $f(z) = \Succ(\Zero(z))$۔ اوس
چې~$h'$ لرو، بيا ترکيب کارولے شو او~$h(y) =
h'(\Proj{1}{0}(y),\Proj{1}{0}(y))$ ټاکلے شو۔ دا ښيي چې~$h$ له
بنسټيزو تابعو څخه د ترکيبونو او بنسټيزو بازګښتونو د يوې لړۍ په
کارولو تعريفېدے شي، نو~$h$ بنسټيزه بازګشتي ده۔
\end{ex}

\end{document}
